\hypertarget{chapter-22-the-half-the-world}{%
\section{Chapter 22: The Half the
World}\label{chapter-22-the-half-the-world}}

\begin{quote}
\itshape ``Give an attacker half the world and a lesser chain dies. Give it to
ours and all it earns is the right to be seen lying.''\\
\normalfont\hfill---DagKnight consensus notes
\end{quote}

\hypertarget{scene-1-the-majority-that-wasnt-enough}{%
\subsection{Scene 1: The Majority That Wasn't
Enough}\label{scene-1-the-majority-that-wasnt-enough}}

The coalition came back, and this time it did not bring a clever exception or a
reasonable tolerance. This time it brought \emph{power}---raw, vulgar, measured in
machines.

``They've bought half the network,'' Sable said, and for once there was no anger in
her, only a kind of awe at the size of it. ``Not half the people. Half the muscle.
Half the hashing, half the witnesses they could lean on or build. The classic
nightmare. On the old chains, that was checkmate---hold half the world's strength
and you decide what history says.''

Elena Voss had read the old textbooks, back when she was learning what she was
defending. A majority of the power meant a majority of the truth: you could outbuild
the honest chain, present your version as the heaviest, and the network, dutifully
following weight, would adopt your lie and call it consensus. It was the oldest
catastrophe in the book, and the powerful had finally simply \emph{bought} it.

``So we've lost,'' a young witness said.

``On a chain that picks a single winner by weight,'' Elena said, ``we would have.
That's why this one never did.''

\hypertarget{scene-2-the-forest-not-the-spine}{%
\subsection{Scene 2: The Forest, Not the
Spine}\label{scene-2-the-forest-not-the-spine}}

The Sigil's consensus was not a single line of blocks marching one after another,
each chosen as the heaviest. It was a \emph{forest}---a structure where many blocks
could grow at once, each pointing back to all the parents it had seen, weaving a
braid rather than a chain. Disagreement was not a fork to be resolved by force; it
was a branch to be absorbed by the structure. Honesty converged not because a
strongest node hammered everyone into line, but because the four committed roots
made every divergence \emph{visible}, and honest nodes simply kept building on
everything they could verify.

``Here's what the coalition never understood,'' Wren said, surfacing again, drawn by
the gravity of the assault. ``On a single-winner chain, owning half the power lets
you \emph{hide} your fork until it's heavier than ours, then reveal it and erase
us. But you can't hide in a forest. The moment their blocks point at a history that
contradicts the committed roots, every honest node \emph{sees} it---not later, not
after the reveal. Immediately. Their half-the-world doesn't buy them a secret. It
buys them a very expensive, very public confession.''

Elena watched the coalition's blocks arrive, heavy and many, and watched the honest
network simply\ldots{} note them. Weigh them. And decline to follow them anywhere
the roots said was false. The attackers had brought a battering ram to a building
with no single door.

\hypertarget{scene-3-the-confession}{%
\subsection{Scene 3: The
Confession}\label{scene-3-the-confession}}

For a few hours it was genuinely frightening, because half the world's power makes
a great deal of noise. The coalition flooded the network with their version of
history, a vast heavy branch in which certain transactions had never happened and
certain balances had always been theirs.

But every block they produced had to commit the four roots, and the four roots do
not lie for anyone, not even for half the world. Their branch was internally
consistent only with \emph{itself}---and the instant it touched the honest braid,
the contradiction lit up across every verifying node on earth. Not a subtle bend
this time. A glaring, root-level, ten-millisecond-checkable lie, broadcast by the
attackers themselves, signed with their own keys, paid for with their own
electricity.

``They spent a fortune,'' Sable breathed, ``to publish a signed admission of
exactly what they tried to steal.''

The honest network did the only thing it ever did. It kept building on the truth it
could verify, and treated the heavy lying branch as what the math said it was: a
large, expensive, self-incriminating dead end. The coalition's half-the-world could
make blocks. It could not make those blocks \emph{agree with reality}, and it could
not make the disagreement \emph{invisible}, and in a forest that watches itself,
visible disagreement is not power. It is evidence.

\hypertarget{scene-4-the-price-of-being-seen}{%
\subsection{Scene 4: The Price of Being
Seen}\label{scene-4-the-price-of-being-seen}}

The coalition withdrew, as the powerful always withdrew once their advantage became
a liability. They had learned, at enormous cost, the lesson the Sigil taught
everyone eventually: that strength is only an advantage in a system where strength
can hide. Take away the hiding, and raw power becomes raw exposure.

``It held,'' Sable said, and then, catching herself, smiled the tired smile.
``It held even against half the world. I genuinely didn't know if it would.''

``It held because we never let truth be a popularity contest,'' Elena said. ``The
old chains said: whatever the most power agrees on is real. We said: whatever the
math can verify is real, and the most power in the world can't make a contradiction
verify. They could outvote us. They could never out-\emph{prove} us. Those are
different countries, and we made sure to live in the second one.''

She thought, as she so often did at these afters, of Berlin---of a girl who had
believed power was the only real currency, that the watchers always won in the end
because they were strong. The watchers had just thrown half the world's strength at
a forest of committed truth, and all their strength had bought them was the
humiliation of being seen.

A stranger on a train verified the tip of a chain that half the world had failed to
rewrite, and looked back out the window, never knowing how close the old nightmare
had come, or how completely the forest had refused it.

Elena raised her hand.

The vigil continued---heavier now, tested against the worst the textbooks
predicted, and still, impossibly, standing.

\hypertarget{chapter-22-notes}{%
\subsection{Chapter Notes}\label{chapter-22-notes}}

\begin{itemize}
\tightlist
\item
  \textbf{The 50\%+ adversary.} On a single-chain, heaviest-wins protocol, an
  attacker controlling a majority of the consensus power can secretly build a
  competing history, then reveal it to overwrite honest blocks---the classic
  majority attack.
\item
  \textbf{DAG consensus changes the game.} A directed-acyclic-graph (``forest'')
  consensus lets many blocks coexist, referencing all parents seen, so divergence
  is absorbed structurally rather than resolved by a single winner. There is no
  hidden fork to reveal.
\item
  \textbf{Committed roots make lies visible.} Because every block commits the four
  state roots, an attacker's branch can be internally consistent only with itself;
  the moment it contradicts the honest braid, the lie is publicly,
  instantly verifiable---broadcast by the attacker's own signed blocks.
\item
  \textbf{Truth is not a popularity contest.} The thematic payoff: the Sigil
  decides reality by what can be \emph{verified}, not by what the most power
  agrees on. Majority power can outvote but cannot out-prove; removing the ability
  to hide turns raw strength into raw exposure.
\end{itemize}
